\documentclass[12pt]{article}
\usepackage[utf8]{inputenc}
\usepackage[T1]{fontenc}
\usepackage{hyperref}
\title{Iberian Language Resources for K-12 Curriculum Development}
\author{Compiled by Claude Sonnet 5.5 (claude-sonnet-5-5) for Andr\'as Kornai}
\date{}
\begin{document}
\maketitle

\section{Introduction}
Iberian (ISO 639-3 xib; Glottocode iber1250) was a Palaeohispanic language of the eastern and southern Iberian Peninsula. It was written in its own scripts and is extinct. The corpus runs to roughly 2,300 inscriptions. The language is still largely undeciphered, so the sounds of the writing can be read but most meanings are unknown. There are no native speakers and I found no school revitalization programme.

\section{Primary Educational Resources}
I found no primers, textbooks or school materials for Iberian. Nothing aimed at pupils turned up.

\section{Linguistic and Reference Materials}
\begin{itemize}
\item Digital Classicist wiki entry on the Hesperia Palaeohispanic Languages Data Bank (corpus of all Palaeohispanic texts): \url{https://wiki.digitalclassicist.org/Hesperia_Palaeohispanic_Languages_Data_Bank}
\item Reichert Verlag, \emph{Lexikon der iberischen Inschriften / L\'exico de las inscripciones ib\'ericas}: \url{https://reichert-verlag.de/en/subjects/linguistics/9783954902910_lexikon_der_iberischen_inschriften_lexico_de_las_inscripciones_ibericas-detail}
\item ``On Hesperia and AELAW'' (paper on the databases): \url{https://latinnow.eu/wp-content/uploads/2019/08/on-hesperia-and-aelaw.pdf}
\item CORDIS project report mentioning the AELAW Iberian volume by Velaza and Gorrochategui: \url{https://cordis.europa.eu/project/id/655938/reporting/pl}
\end{itemize}
A forthcoming \emph{W\"orterbuch der iberischen Inschriften} (Monumenta Linguarum Hispanicarum V.2) is mentioned in the CORDIS report.

\section{Religious and Community Resources}
I found no scripture or community organisations for Iberian. Inscriptions include votive and funerary texts, but I did not see a source for that in these searches.

\section{Digital and Audio Resources}
The Hesperia data bank is the main digital resource. Its project site was not reached directly in my search, so I cite only the Digital Classicist description. I found no audio or video.

\section{Curriculum Development Context}
Material is sparse and mostly scholarly. A K-12 unit could use the Iberian scripts and inscriptions as cultural-heritage content in Spanish, Catalan or English history classes. Because the language is largely undeciphered, a speaking curriculum is not possible.

\end{document}